\documentclass[12pt,a4paper]{article}

\usepackage[utf8]{inputenc}
\usepackage[T1]{fontenc}
\usepackage[french]{babel}
\usepackage{geometry}
\geometry{margin=2.5cm}
\usepackage{hyperref}
\usepackage{xcolor}
\usepackage{enumitem}
\usepackage{parskip}
\usepackage{listings}

\lstdefinestyle{bash}{
  basicstyle=\ttfamily\small,
  keywordstyle=\color{black},
  commentstyle=\color{gray},
  stringstyle=\color{teal},
  showstringspaces=false,
  breaklines=true,
  frame=single,
  rulecolor=\color{black!15},
  backgroundcolor=\color{black!2},
  numbers=none
}

\title{Rapport de projet (PFA Frontend) --- Application de gestion du pointage numérique}
\author{}
\date{}

\begin{document}
\maketitle
\tableofcontents
\newpage

\textbf{Projet} : TimeTrack Pro --- Gestion du pointage numérique des employés\\
\textbf{Référentiel} : Cahier des charges PFA + projet Frontend (« Application de gestion du pointage numérique des employés »)\\
\textbf{Technologies livrées} : Next.js (App Router), Firebase (Auth + Firestore), Firebase Cloud Functions, Tailwind CSS, Recharts

\begin{quote}
Ce rapport décrit le produit réalisé, son architecture, ses fonctionnalités, ainsi que la conformité au cahier des charges et les améliorations apportées.
\end{quote}

\section{Contexte et objectif}
L’objectif du projet est de développer une application web permettant de :
\begin{itemize}[leftmargin=*]
  \item gérer le \textbf{pointage} des employés en temps réel,
  \item vérifier la présence via \textbf{géolocalisation (GPS)},
  \item scanner un \textbf{QR code sécurisé},
  \item gérer les \textbf{congés},
  \item offrir un \textbf{tableau de bord analytique},
  \item détecter les anomalies : \textbf{retards}, \textbf{absences}, \textbf{sorties anticipées} et \textbf{insuffisances d’heures}.
\end{itemize}

\section{Périmètre et acteurs}
\subsection{Acteurs}
\begin{itemize}[leftmargin=*]
  \item \textbf{Administrateur} : pilote l’ensemble (dashboard, employés, congés, rapports).
  \item \textbf{Employé} : pointe (GPS + QR), consulte son historique, demande des congés, gère son profil.
\end{itemize}

\subsection{Hypothèses et périmètre}
\begin{itemize}[leftmargin=*]
  \item L’application est pensée pour un usage \textbf{web} (desktop + mobile).
  \item La validation du pointage (QR + zone) se fait \textbf{côté backend}.
  \item Les données opérationnelles sont stockées dans \textbf{Firestore}.
\end{itemize}

\section{Architecture technique (réalisée)}
\subsection{Frontend}
\begin{itemize}[leftmargin=*]
  \item \textbf{Framework} : Next.js \texttt{16.x} (App Router).
  \item \textbf{UI \& UX} : Tailwind CSS + composants UI (Radix).
  \item \textbf{Formulaires \& validation} : React Hook Form + Zod.
  \item \textbf{Graphiques} : Recharts.
  \item \textbf{Notifications} : Sonner (toasts).
\end{itemize}

\subsection{Backend}
\begin{itemize}[leftmargin=*]
  \item \textbf{Firebase Cloud Functions (v2)} : logique serveur pour le pointage (Callable Function).
  \item \textbf{Auth Trigger (recommandation/évolution)} : création automatique du profil \texttt{users/\{uid\}} à la création d’un compte (voir \S 11).
\end{itemize}

\subsection{Données}
\begin{itemize}[leftmargin=*]
  \item \textbf{Stockage principal} : Firestore (\texttt{users}, \texttt{pointages}, \texttt{conges}).
\end{itemize}

\begin{quote}
\textbf{Remarque conformité CDC :} le cahier des charges propose MongoDB/REST en exemple, mais n’impose pas strictement la technologie. Le résultat fonctionnel (collections, sécurité, API) est équivalent. La solution Firestore + Functions garantit la sécurité et simplifie la mise en œuvre.
\end{quote}

\section{Modèle de données (Firestore)}
Le modèle suit l’esprit des collections attendues :

\subsection{\texttt{users}}
\begin{itemize}[leftmargin=*]
  \item \texttt{nom}, \texttt{email}, \texttt{role} (\texttt{admin} \,|\, \texttt{employe}), \texttt{createdAt}
  \item champs profil employé (optionnels) : \texttt{matricule}, \texttt{telephone}, \texttt{departement}, \texttt{poste}, \texttt{cin}, \texttt{adresse}, \texttt{dateNaissance}, \texttt{dateEmbauche}
\end{itemize}

\subsection{\texttt{pointages}}
\begin{itemize}[leftmargin=*]
  \item \texttt{userId}, \texttt{date} (YYYY-MM-DD), \texttt{heure} (HH:MM), \texttt{type} (\texttt{entree} \,|\, \texttt{sortie})
  \item \texttt{latitude}, \texttt{longitude}, \texttt{valide}, \texttt{createdAt} (server timestamp)
  \item (optionnel) précision GPS \texttt{accuracyM}
\end{itemize}

\subsection{\texttt{conges}}
\begin{itemize}[leftmargin=*]
  \item \texttt{userId}, \texttt{dateDebut}, \texttt{dateFin}, \texttt{type} (\texttt{annuel} \,|\, \texttt{maladie} \,|\, \texttt{exceptionnel})
  \item \texttt{statut} (\texttt{en\_attente} \,|\, \texttt{valide} \,|\, \texttt{refuse}), \texttt{createdAt}
\end{itemize}

\section{Fonctionnalités principales (réalisées)}
\subsection{Authentification \& rôles}
Conforme à la demande :
\begin{itemize}[leftmargin=*]
  \item \textbf{Connexion sécurisée} email + mot de passe.
  \item \textbf{Rôles} admin / employé.
  \item Gestion de session côté client via Firebase Auth + lecture du rôle dans Firestore.
\end{itemize}

\subsection{Pointage intelligent (GPS + QR)}
Flux employé (conforme CDC) :
\begin{enumerate}[leftmargin=*]
  \item L’employé ouvre la page \textbf{Pointer}.
  \item L’application récupère la \textbf{géolocalisation} (stratégie en 2 étapes : rapide puis haute précision).
  \item L’employé scanne un \textbf{QR code}.
  \item Le backend valide :
  \begin{itemize}
    \item token QR attendu,
    \item présence dans la \textbf{zone autorisée},
    \item puis enregistre le pointage (entrée/sortie) avec timestamp serveur.
  \end{itemize}
\end{enumerate}

Améliorations UX / anti-fraude :
\begin{itemize}[leftmargin=*]
  \item feedback clair (état GPS, précision, distance, zone OK/KO),
  \item déduplication des notifications « QR détecté »,
  \item stabilisation de la caméra (éviter flux doublé + AbortError).
\end{itemize}

\subsection{Gestion des congés}
\textbf{Employé} :
\begin{itemize}[leftmargin=*]
  \item demande de congé avec type, dates, et calcul de durée (jours ouvrés).
\end{itemize}

\textbf{Admin} :
\begin{itemize}[leftmargin=*]
  \item visualisation des demandes,
  \item \textbf{validation / refus},
  \item filtres (en attente / tous).
\end{itemize}

\subsection{Tableau de bord analytique (admin)}
Dashboard admin :
\begin{itemize}[leftmargin=*]
  \item compteurs (utilisateurs, pointages, congés),
  \item graphique pointages (7 jours),
  \item distribution des congés par statut (pie chart),
  \item design modernisé (cartes, gradients).
\end{itemize}

\subsection{Filtres avancés et rapports}
\begin{itemize}[leftmargin=*]
  \item Historique employé : filtres (type, mois), KPIs (retards/absences/présence), export \textbf{CSV}.
  \item Liste employés admin : recherche multi-champs + tri + export \textbf{CSV}.
  \item Rapport pointage admin : filtres + export \textbf{CSV}.
\end{itemize}

\subsection{Détection d’anomalies (MVP + base)}
Le projet inclut une première détection visuelle côté historique (retards/anomalies).\\
Les règles exactes « heures < seuil 8h/jour » sont prévues en extension (voir \S 11).

\section{Interfaces utilisateur (réalisées)}
\subsection{Espace Employé}
\begin{itemize}[leftmargin=*]
  \item Connexion
  \item Pointage (GPS + QR)
  \item Historique personnel + KPIs + export
  \item Demande de congé
  \item Profil employé (édition)
  \item Navigation responsive (hamburger menu)
  \item \textbf{Light/Dark mode}
\end{itemize}

\subsection{Espace Admin}
\begin{itemize}[leftmargin=*]
  \item Dashboard analytique (charts)
  \item Liste employés (profil complet + tri)
  \item Gestion congés (valider/refuser)
  \item Rapport pointage (export)
  \item Navigation responsive (hamburger menu)
  \item \textbf{Light/Dark mode}
\end{itemize}

\section{Sécurité (réalisée)}
Conforme aux attentes :
\begin{itemize}[leftmargin=*]
  \item \textbf{Authentification requise} pour les zones applicatives.
  \item Validation backend obligatoire pour le \textbf{pointage} :
  \begin{itemize}
    \item QR token,
    \item zone GPS.
  \end{itemize}
  \item Contrôle d’accès par \textbf{Firestore Rules} (lecture utilisateur, actions admin).
\end{itemize}

Points clés :
\begin{itemize}[leftmargin=*]
  \item Le rôle admin/employé est stocké en base et utilisé pour autoriser/filtrer les pages.
  \item Le QR code peut être « rotatif » via simple changement de token côté configuration (paramètres functions).
\end{itemize}

\section{Contraintes techniques (réalisées)}
\begin{itemize}[leftmargin=*]
  \item \textbf{Responsive design} (mobile obligatoire) : réalisé avec menus mobiles.
  \item Performance / UX :
  \begin{itemize}
    \item géolocalisation optimisée,
    \item scans et toasts stabilisés,
    \item requêtes Firestore limitées (ex. 7 jours via \texttt{where(in)}).
  \end{itemize}
\end{itemize}

Temps réel / WebSocket :
\begin{itemize}[leftmargin=*]
  \item présent comme piste (socket.io dépendance), mais non requis pour le MVP.
\end{itemize}

\section{Tests et validation}
\subsection{Tests fonctionnels (manuel)}
\begin{itemize}[leftmargin=*]
  \item Création compte + connexion.
  \item Pointage : GPS OK/KO + QR OK/KO.
  \item Congés : création employé, validation/refus admin.
  \item Dashboard : KPI + graphiques.
  \item Exports CSV : admin et employé.
  \item Responsive : mobile + hamburger menu.
  \item Thème : light/dark toggle.
\end{itemize}

\subsection{Qualité}
\begin{itemize}[leftmargin=*]
  \item ESLint intégré et vérifié.
\end{itemize}

\section{Déploiement \& exécution}
\subsection{Lancement local (frontend)}
\begin{lstlisting}[style=bash]
npm install
npm run dev
\end{lstlisting}

\subsection{Fonctions (émulateur)}
\begin{lstlisting}[style=bash]
cd functions
npm install
npm run serve
\end{lstlisting}

\begin{quote}
Note : le déploiement Functions v2 peut nécessiter un plan Firebase adapté (Blaze) selon le contexte. L’émulateur reste la solution locale recommandée.
\end{quote}

\section{Évolutions et recommandations (professionnelles)}
\subsection{Simplification “pro” de l’inscription (recommandée)}
Pour supprimer les problèmes de synchronisation rôle/profil après inscription :
\begin{itemize}[leftmargin=*]
  \item ajouter une Cloud Function \textbf{Auth Trigger} \texttt{onCreate} qui crée systématiquement \texttt{users/\{uid\}},
  \item ne plus écrire \texttt{users/\{uid\}} depuis le client,
  \item simplifier les Firestore Rules (moins permissives).
\end{itemize}

\subsection{Détection d’anomalies avancée}
\begin{itemize}[leftmargin=*]
  \item agrégation « heures/jour < 8h »,
  \item retards configurables,
  \item sorties anticipées basées sur planning.
\end{itemize}

\subsection{Exports}
\begin{itemize}[leftmargin=*]
  \item Export \textbf{PDF / Excel} côté admin.
\end{itemize}

\subsection{Renforcement sécurité}
\begin{itemize}[leftmargin=*]
  \item rotation programmée du QR token,
  \item logs d’audit (admin actions),
  \item rate limiting (anti-spam pointage).
\end{itemize}

\section{Conclusion}
Le projet TimeTrack Pro fournit une application web moderne répondant au cahier des charges :
\begin{itemize}[leftmargin=*]
  \item pointage intelligent (GPS + QR + validation backend),
  \item gestion des congés employé/admin,
  \item dashboard analytique avec graphiques,
  \item filtres, exports, détection d’anomalies (MVP),
  \item interface responsive et thème light/dark,
  \item architecture sécurisée et évolutive.
\end{itemize}

\section*{Annexes (à compléter si demandé)}
\begin{itemize}[leftmargin=*]
  \item Diagrammes UML (cas d’utilisation, séquence pointage, MCD/collections).
  \item Captures d’écran des interfaces.
  \item Démo vidéo / lien repo GitHub.
\end{itemize}

\end{document}

